\begin{resumo}
Modelos de fatores são amplamente empregados na precificação de ativos em finanças quantitativas. Eles decompõem os retornos em exposições a um conjunto reduzido de variáveis comuns, concentrando a informação relevante do mercado sem impor estrutura econômica prévia. Avanços recentes em aprendizado de máquina permitiram que esses fatores fossem extraídos diretamente dos dados, tratando-os como variáveis latentes aprendidas conjuntamente ao processo de previsão. O \textit{FactorVAE} estende essa abordagem ao modelar os fatores latentes como variáveis aleatórias dentro de um \textit{autoencoder} variacional, o que permite representar explicitamente a incerteza associada ao sinal extraído, propriedade relevante dada a baixa relação sinal-ruído característica de séries financeiras. Este trabalho replica e avalia o \textit{FactorVAE} no mercado brasileiro de ações, adaptando o arcabouço original ao universo de ativos da B3. A questão central é se a vantagem de ordenação cross-sectional que o modelo apresenta no mercado chinês se preserva em um ambiente com menor liquidez e menor número de ativos negociados. O desempenho é avaliado em relação a três baselines (GRU, CA, IPCA), e os resultados são analisados tanto em termos de capacidade preditiva quanto de aplicabilidade na construção de carteiras.

\textbf{Palavras-chaves}: \imprimirpalavraschavesresumo.
\end{resumo}